\begin{proposition}
    基本集合 $A$ に対し，
    \[
    m_J(A) = \lambda^*(A).
    \]
\end{proposition}

\begin{proof}
\[
\lambda^*(A)\le m_J(A)
\]
は，$A$ 自身による有限被覆を考えれば従う．

逆向きの不等式を示すため，
$A \subset \bigcup_{n=1}^{\infty} R_n$ を満たす任意の可算被覆
$\{R_n\}_{n=1}^{\infty} \subset \calE_d$ を取る．
もし
\[
\sum_{n=1}^{\infty}|R_n|=\infty
\]
なら
\[
m_J(A)\le \sum_{n=1}^{\infty}|R_n|
\]
は自明である．
したがって右辺が有限であるとする．このとき各 $R_n$ は有界である．

$\eps>0$ を取り，コンパクトなJordan可測集合 $K\subset A$ を任意に固定する．
各 $n$ について $R_n$ を少し膨らませた開区間 $U_n$ を
\[
R_n \subset U_n,
\qquad
|U_n| \le |R_n|+\frac{\eps}{2^n}
\]
となるように取る．
すると $\{U_n\}_{n=1}^{\infty}$ はコンパクト集合 $K$ の開被覆であるから，
有限部分被覆
\[
K \subset U_{n_1}\cup\cdots\cup U_{n_N}
\]
が存在する．
Jordan外測度の有限劣加法性より
\[
m_J(K)=m_J^*(K)
\le
\sum_{j=1}^N m_J^*(U_{n_j})
\le
\sum_{j=1}^N |U_{n_j}|
\le
\sum_{n=1}^{\infty}|R_n|+\eps
\]
である．
したがって $\eps>0$ は任意なので
\[
m_J(K)\le \sum_{n=1}^{\infty}|R_n|
\]
である．
一方，基本集合 $A$ を互いに素な半開区間の合併
\[
A=\bigsqcup_{j=1}^N Q_j,
\qquad
Q_j=\prod_{i=1}^d(a_{ij},b_{ij}]
\]
と書く．
したがって，$m_J(A)<\infty$ のときは，
各 $Q_j$ の端点は有限であり，
\[
K_\eta
:=
\bigsqcup_{j=1}^N \prod_{i=1}^d [a_{ij}+\eta,b_{ij}]
\subset A
\]
は十分小さい $\eta>0$ に対してコンパクトなJordan可測集合である．
また $m_J(K_\eta)\to m_J(A)$ $(\eta\downarrow0)$ であるから，
任意の $\delta>0$ に対して $\eta>0$ を十分小さく取れば，
あるコンパクトなJordan可測集合 $K_\delta\subset A$ が
\[
m_J(A)-\delta \le m_J(K_\delta)
\]
を満たす．
これを上の不等式に代入すると
\[
m_J(A)-\delta
\le
m_J(K_\delta)
\le
\sum_{n=1}^{\infty}|R_n|
\]
である．
$\delta>0$ は任意なので
\[
m_J(A)\le \sum_{n=1}^{\infty}|R_n|
\]
である．

また，$m_J(A)=\infty$ のときは，
ある $Q_j$ の体積が無限大である．
この $Q_j$ の無限端点を有限の端点で切り詰め，
有限端点は必要なら少し内側へ動かすことで，
任意の $M>0$ に対してコンパクトなJordan可測集合
$K_M\subset Q_j\subset A$ を $m_J(K_M)\ge M$ となるように取れる．
すると上の不等式から
\[
M\le \sum_{n=1}^{\infty}|R_n|
\]
が任意の $M>0$ について成り立つことになり，
これはいま仮定している $\sum_{n=1}^{\infty}|R_n|<\infty$ に矛盾する．
したがって有限和の可算被覆を考える場合にも
\[
m_J(A)\le \sum_{n=1}^{\infty}|R_n|
\]
である．
被覆について下限を取れば
\[
m_J(A)\le \lambda^*(A)
\]
である．
\end{proof}

% \begin{remark}
% ここではJordan外測度の可算劣加法性は使っていない．
% 実際，Jordan外測度は一般には可算劣加法的ではない．
% 例えば
% \[
% \QQ\cap[0,1]
% =
% \bigcup_{q\in \QQ\cap[0,1]}\{q\}
% \]
% であり，各一点集合のJordan外測度は $0$ であるが，
% $\QQ\cap[0,1]$ のJordan外測度は $1$ である．
% \end{remark}

\begin{lemma}[Jordan可測集合の基本集合近似]
Jordan可測集合 $A$ と $\eps>0$ に対し，
基本集合 $F,G \in \calA_d$ が存在して
\[
F \subset A \subset G,
\qquad
m_J(G)-m_J(F)<\eps
\]
を満たす．
\end{lemma}

\begin{proof}
\red{第2章のJordan可測性の近似定理}より，
基本集合 $F,G \in \calA_d$ が存在して
\[
F \subset A \subset G,
\qquad
m_J(G\setminus F)<\eps
\]
を満たす．
基本集合の有限加法性より
\[
m_J(G)=m_J(F)+m_J(G\setminus F)
\]
だから
\[
m_J(G)-m_J(F)<\eps
\]
である．
\end{proof}

\begin{theorem}
    Jordan可測集合 $A$ に対し，
    \[
    m_J(A) = \lambda^*(A).
    \]
\end{theorem}

\begin{proof}
Jordan可測集合 $A$ は有界である．
任意の $\eps > 0$ に対し，前補題より
基本集合 $F,G \in \calA_d$ を
\[
F \subset A \subset G,
\qquad
m_J(G)-m_J(F)<\eps
\]
となるように取る．
単調性と前命題より
\[
\lambda^*(F)=m_J(F)\le \lambda^*(A)\le \lambda^*(G)=m_J(G)
\]
である．
また $F \subset A \subset G$ だから
\[
m_J(F)\le m_J(A)\le m_J(G)
\]
であり，
\[
m_J(G)-m_J(F)<\eps
\]
だから
\[
m_J(A)-\eps < m_J(F)=\lambda^*(F)
\le
\lambda^*(A)
\le
\lambda^*(G)=m_J(G)
< m_J(A)+\eps
\]
を得る．
$\eps > 0$ は任意なので $\lambda^*(A)=m_J(A)$ である．
\end{proof}
